\documentclass[10pt]{article}

\usepackage[margin=1in]{geometry}
\usepackage{amsmath,amssymb}
\usepackage{booktabs}
\usepackage[round]{natbib}
\usepackage{microtype}

\newcommand{\one}{\mathbf{1}}

\title{Bequest specification for a housing--fertility lifecycle model:\\
what the motive should depend on}
\author{Fertility--Housing project}
\date{14 July 2026}

\begin{document}
\maketitle

\section{The question}

The object at issue is a primitive: the utility a household derives from the
estate it leaves at death, as a function of the estate size and of the
household's completed number of children $n$. The active model gives this
terminal payoff the form
%
\begin{equation}
v(b,n)=\theta_0\,\max\{1+\theta_n n,\,0\}
\left[\frac{(\theta_1+b)^{1-\sigma}}{1-\sigma}
-\frac{\theta_1^{1-\sigma}}{1-\sigma}\right],
\qquad \theta_1=0.01,\ \sigma=2,
\label{eq:current}
\end{equation}
%
where $b$ is the gross estate realized at certain death at age~82 and every
household is endowed with the motive.\footnote{Form transcribed from
\texttt{code/model/intergen\_housing\_fertility/solver.py}
(\texttt{bequest\_utility\_vec}); not modified here.} Two features of
\eqref{eq:current} do economic work: the multiplier $\max\{1+\theta_n n,0\}$
makes the \emph{strength} of the motive rise linearly in the number of
children, and the near-zero shifter $\theta_1$ makes the term a nearly
homothetic warm glow rather than a luxury good that only the rich act on.

The design question has an extensive and a functional part. Extensively, is
the motive a property of parents, so that childless households leave estates
only accidentally? Functionally, if the motive is present, does its strength
rise with the number of children, and is a linear estate multiplier the right
way to encode that? The two parts have different answers, and neither answer
is the current specification.

\section{Evidence}

\paragraph{(a) The extensive margin of the motive.} The best direct evidence
finds at most an extensive shifter, not a strong gradient in $n$.
\citet{kopczuk2007} estimate that about three-quarters of elderly singles hold
a bequest motive; children raise the probability of a motive from $63\%$ to
$79\%$, significant only at the $10\%$ level, and the remaining child
covariates (income, education, proximity) are mostly insignificant. They read
the evidence as an egoistic motive but explicitly reject the older assumption,
due to \citet{hurd1989}, that only parents have a motive at all---an
assumption \citeauthor{hurd1989} used to rationalize faster parental
decumulation as accidental bequests. The margin the data support is that
parents are more likely to \emph{have} a motive, not that the motive scales
with how many children they have.

\paragraph{(b) Children matter, but through transfers and division.} The
operative child channel in the modern evidence is inter vivos giving, not the
shape of the terminal estate utility. \citet{kvaerner2023} uses quasi-random
parental cancer diagnoses in Norwegian administrative data: a diagnosis raises
children's wealth by roughly $10\%$ while dynasty wealth is unchanged, which
identifies deliberate, child-directed \emph{transfers} and re-reads
\citeauthor{hurd1989}'s faster decumulation as giving rather than
indifference. Against a strong altruistic reading, \citet{altonji1997} find
that redistributing one dollar from parent to child moves the transfer by less
than thirteen cents, rejecting the sharp Barro--Becker benchmark.
\citet{mcgarry1999} shows that inter vivos transfers are compensatory---they
flow to lower-income children---yet bequests are divided \emph{equally} in
$83\%$ of wills, and \citet{wilhelm1996} finds $88\%$ equal division in estate
records, regardless of child income. The terminal estate is therefore close to
a mechanical equal split; the child-contingent action is in the timing and
targeting of transfers before death.

\paragraph{(c) Accumulation by asset class.} Whether parents ``save more''
depends entirely on the asset, and the pattern favors housing as the
child-directed stock. \citet{scholzseshadri} find that children have a large
negative effect on financial net worth over working life: families with
children consume more and optimally hold less retirement wealth, a crowd-out
rather than an absence of caring. \citet{boar2021} shows the mirror image in
consumption: parental consumption responds to the \emph{child's} permanent
income risk with an elasticity of $-0.06$ to $-0.08$, precautionary saving
held as insurance for the child. \citet{bernheim1991} infers bequest motives
from the life-insurance and annuity choices of the elderly, a further
child-contingent protective margin. Housing runs the other way: the housing
budget share rises with the number of children
\citep{coeurdacier2023}, consistent with this repository's own PSID
birth event studies. The synthesis is that ``parents save more'' is true in
houses, insurance, transfers, and buffers and false in measured financial
wealth---and housing is exactly the child-directed asset the model already
carries.

\paragraph{(d) Structural conventions.} The structural literature encodes the
motive as a child-blind luxury warm glow. \citet{denardi2004} introduces the
luxury shifter that makes bequests a good the rich act on;
\citet{lockwood2018} and \citet{ameriks2011} estimate the same
$v(X)=\frac{a}{1-\sigma}\big(X/a+\phi\big)^{1-\sigma}$ form, with $\phi$ the
threshold below which nothing is left. \citet{denardi2025} find bequests
behave as luxuries concentrated among rich couples, with children not even a
conditioning variable in the saving heterogeneity. Closest to this project,
\citet{nakajima2017} estimate a child-blind warm glow inside a housing model,
$\gamma\,(a+\zeta)^{1-\sigma}/(1-\sigma)$, with strength $\gamma=0.43$ and
shifter $\zeta\approx\$19{,}600$---a luxury threshold of order a quarter of
annual household income. Even \citet{kvaerner2023}, the paper with the
strongest child-directed evidence, models the motive itself child-blind and
puts children in the transfer block.\footnote{Verification status, following
the round-3 source memo's tags. Independently verified today [V]: the negative
sign of children on financial net worth (\citeauthor{scholzseshadri}), the
$-0.06$ to $-0.08$ child-risk consumption elasticity (\citeauthor{boar2021}),
the rising housing budget share in $n$ (\citeauthor{coeurdacier2023}), and the
life-insurance reading (\citeauthor{bernheim1991}, at search level).
Extracted from paper pages or abstracts today but not yet re-checked against
the papers [E], and load-bearing: \citeauthor{kopczuk2007}'s $63\%\!\to\!79\%$
and $10\%$ significance, \citeauthor{nakajima2017}'s $\gamma=0.43$ and
$\zeta\approx\$19{,}600$, \citeauthor{kvaerner2023}'s $\sim\!10\%$ wealth
effect, and \citeauthor{mcgarry1999}'s $83\%$ / \citeauthor{wilhelm1996}'s
$88\%$ equal-division shares. Spot-check these before paper use.}

\section{Model evidence}

The current block fails for a reason internal to the model, and the failure is
sharp. With certain death at age~82, no survival risk, and no downsizing
force, the cheapest way to hold an estate is housing equity, so the motive
resolves into a single instruction: die owning the largest possible house,
leveraged. Table~\ref{tab:mech} shows what switching pure warm glow on does at
the clean frontier coordinates.

\begin{table}[h]
\centering
\small
\begin{tabular}{lccc}
\toprule
object & $\theta_0=0$ & $\theta_0=0.3$ & data target \\
\midrule
ownership at ages $74+$ & $0.606$ & $0.981$ & $\sim\!0.76$ \\
mean liquid wealth $74\!+\!/\,62\text{--}74$ & $+0.17$ & $-0.33$ & --- \\
old nonhousing wealth median & $2.33$ & $3.30$ & $2.23$ \\
TFR & $2.004$ & $2.000$ & $1.918$ \\
\bottomrule
\end{tabular}
\caption{Fixed-$\theta$ profiles, $\theta_n=0$, tight evaluator. Warm glow
drives old-age ownership to $98\%$ and mean liquid wealth at $74+$ negative:
households borrow against the house, consume, and bequeath the house.}
\label{tab:mech}
\end{table}

The old-age ownership target ($0.764$, weight $160$) makes this arithmetically
unaffordable, and it is right to bind: the data say old households do not all
keep their houses. Reoptimizing everything else with $\theta_0=0.3$ forced
lands at $9.84$--$10.42$ against a verified no-bequest frontier of $6.96$ tight,
so the reoptimized price of pure warm-glow bequests is about $+2.9$ loss
points; even the smallest tested warm glow, $\theta_0=0.05$, costs $+2.0$. The
linear child multiplier is a second, separable catastrophe, multiplying the
penalty roughly sevenfold on the audit's accounting: at $\theta_0=0.3$ the
loss rises from $11.43$ under pure warm glow to $96.83$ with the multiplier
on, as the estate becomes a fertility subsidy stacked on the housing-retention
channel.\footnote{Diagnostic value $\theta_n=0.375$; the production default is
$\theta_n=0.25$. The parameter is an exact null direction at $\theta_0\approx0$,
since the multiplier scales a term whose level is zero.} Inverse-variance
reweighting does not reverse the verdict: under SE-proportional weights the
no-bequest baseline scores $345.1$ against $407.2$--$447.8$ for the best
bequest chain, because precision weighting loads even harder on the ownership
moments that bequests wreck. Bequests are not rejected by the data or the
targets; they are rejected by an accounting identity---estates can only be
houses---that the model change in Section~4 is designed to break.

\section{Proposed specification}

\paragraph{Baseline.} Replace \eqref{eq:current} with a child-blind luxury
warm glow that parents, and only parents, carry:
%
\begin{equation}
v(b,n)=\one\{n\ge 1\}\,\theta_0\,
\frac{(\theta_1+b)^{1-\sigma}}{1-\sigma},
\qquad \theta_1\ \text{estimated.}
\label{eq:baseline}
\end{equation}
%
The indicator encodes the one margin the data support---parents are more
likely to have a motive \citep{kopczuk2007,hurd1989}---while the strength is
child-blind in line with \citet{denardi2004,denardi2025,nakajima2017}, and
$\theta_1$ is promoted from $0.01$ to an estimated luxury scale of order a
quarter of annual income, anchored near the \citet{nakajima2017} shifter.
Identification is clean: with the two free bequest parameters
$(\theta_0,\theta_1)$, the old nonhousing wealth median (target $2.23$)
disciplines the luxury threshold and the decumulation moment below disciplines
the strength.

\paragraph{Where the child channel goes.} The intensive dependence on $n$
moves out of the estate utility and into inter vivos down-payment transfers at
family formation, where the causal and descriptive evidence places it
\citep{kvaerner2023,mcgarry1999}, and where the down-payment-gift literature
locates it for constrained first-time buyers.\footnote{Engelhardt and Mayer,
down-payment-gift evidence; the exact year and outlet are not verified from
within the repository and the citation is deliberately left unpinned.} This is
also where a housing--fertility mechanism naturally lives, since the transfer
relaxes the very $(1-\phi)$ down-payment constraint that couples young
ownership to childbearing. Dropping $\theta_n$ costs no identification: it is a
null direction at $\theta_0\approx0$, no literature supports the linear estate
multiplier, and its economic content is reassigned to the transfer block,
which carries its own identifying moment in birth-timed gift incidence.

\paragraph{Equal-division robustness variant.} If a child term in the estate is
wanted at all, the disciplined form comes from equal division rather than a
free slope. Under equal division each of $n$ children receives $b/n$;
aggregating $n$ identical CRRA continuation values with a Barro--Becker
altruism weight $n^{\epsilon}$ \citep{barrobecker1989} gives
%
\begin{equation}
\theta_0\,n^{\epsilon}\,\frac{(\theta_1+b/n)^{1-\sigma}}{1-\sigma}
\ \xrightarrow[\ \theta_1\to0\ ]{}\
\theta_0\,n^{\,\epsilon+\sigma-1}\,\frac{b^{1-\sigma}}{1-\sigma},
\label{eq:equaldiv}
\end{equation}
%
so equal division plus altruism collapses to a warm glow whose child-count
scale is $n^{\epsilon+\sigma-1}$, not a free linear multiplier. For $\sigma=2$
the exponent is $\epsilon+1$, which lies in $[1,2]$ for $\epsilon\in[0,1]$---a
bounded, micro-founded range in place of the unrestricted $\theta_n$; the exact
form $\theta_0\,n\,u(\theta_1+b/n)$ is the $\epsilon=1$ endpoint. This variant
is justified by the equal-division fact \citep{mcgarry1999,wilhelm1996}, not by
a stronger motive for parents, and it should be reported as robustness, never
as the baseline.

\paragraph{A first bounded diagnostic (14 July).} A 53-solve exercise tested
the parent-gated form \eqref{eq:baseline} at
$(\theta_0,\theta_1)\in\{0.05,0.30\}\times\{0.25,1.00\}$, with and without an
owner loan-to-value cap tapering linearly from age~66 to zero at~82, holding
the live 15-moment objective and tight evaluator fixed.\footnote{Packet:
\texttt{output/model/intergen\_bequest\_calibration\_exercise\_20260714/};
summary in the top section of \texttt{CALIBRATION\_STATUS.md}.} The gated form
alone is loss-neutral: the best positive-bequest point scores $6.9774$ against
the no-bequest frontier $6.9647$, so the respecified utility neither helps nor
harms the fit while the exit margin is missing. The taper arm reaches
$4.3718$, but the gain is an artifact of where the data are measured: the
forced schedule generates an ownership cliff---$0.860$ at age~70, $0.561$
at~74, near zero at~78---that matches the 65--75 ownership average ($0.766$
against $0.764$) precisely because the collapse lands just outside the
measured window, while the $74+\,/\,62\text{--}74$ liquid-wealth ratio stays
at $1.016$, so the model still produces no retirement decumulation. Households
are pushed out of houses without being induced to run down wealth, the
opposite of the missing margin. The diagnostic therefore sharpens the
prerequisite: the exit margin must be an empirically pinned mortgage-balance
or downsizing profile, not a forced linear-to-zero loan-to-value schedule, and
the moment set must include $75+$ ownership and a $75+\,/\,62\text{--}74$
decumulation ratio before $\theta_0$ is re-estimated.

\paragraph{Structural prerequisites (model changes awaiting approval).} None of
\eqref{eq:baseline}--\eqref{eq:equaldiv} should be estimated until the
accounting identity of Section~3 is broken. That requires an old-age exit
margin---an externally pinned post-retirement amortization schedule or survival
risk before age~82---so that estates can be financial rather than only
housing, and a $75+\,/\,62\text{--}74$ decumulation moment so that $\theta_0$ is
identified against $\beta$ rather than confounded with it. To be explicit about
scope, the following are proposals, not implemented changes: (i) the old-age
exit margin; (ii) the new decumulation target; (iii) the $\one\{n\ge1\}$
respecification of the motive; (iv) promoting $\theta_1$ to an estimated
parameter; and (v) removing the $(1+\theta_n n)$ multiplier. Items (i) and (ii)
are the load-bearing model changes and should be approved before $\theta_0$ is
re-estimated.

\begingroup
\small
\begin{thebibliography}{99}

\bibitem[Altonji et al.(1997)]{altonji1997}
Altonji, J.~G., F.~Hayashi, and L.~J. Kotlikoff (1997).
Parental altruism and inter vivos transfers: Theory and evidence.
\emph{Journal of Political Economy} 105(6), 1121--1166.

\bibitem[Ameriks et al.(2011)]{ameriks2011}
Ameriks, J., A.~Caplin, S.~Laufer, and S.~Van Nieuwerburgh (2011).
The joy of giving or assisted living? Using strategic surveys to separate
public care aversion from bequest motives.
\emph{Journal of Finance} 66(2), 519--561.

\bibitem[Barro and Becker(1989)]{barrobecker1989}
Barro, R.~J. and G.~S. Becker (1989).
Fertility choice in a model of economic growth.
\emph{Econometrica} 57(2), 481--501.

\bibitem[Bernheim(1991)]{bernheim1991}
Bernheim, B.~D. (1991).
How strong are bequest motives? Evidence based on estimates of the demand for
life insurance and annuities.
\emph{Journal of Political Economy} 99(5), 899--927.

\bibitem[Boar(2021)]{boar2021}
Boar, C. (2021).
Dynastic precautionary savings.
\emph{Review of Economic Studies} 88(6), 2735--2765.

\bibitem[Coeurdacier et al.(2023)]{coeurdacier2023}
Coeurdacier, N., P.-P. Combes, L.~Gobillon, and F.~Oswald (2023).
Fertility, housing costs and city growth. Unpublished manuscript.

\bibitem[De Nardi(2004)]{denardi2004}
De Nardi, M. (2004).
Wealth inequality and intergenerational links.
\emph{Review of Economic Studies} 71(3), 743--768.

\bibitem[De Nardi et al.(2025)]{denardi2025}
De Nardi, M., E.~French, J.~B. Jones, and R.~McGee (2025).
Why do couples and singles save during retirement? Household heterogeneity and
its aggregate implications.
\emph{Journal of Political Economy} (forthcoming).

\bibitem[Engelhardt and Mayer(n.d.)]{engelhardtmayer}
Engelhardt, G.~V. and C.~J. Mayer.
Down-payment gifts and constrained first-time homebuyers.
Down-payment-gift transfer literature; exact year and outlet not verified from
within the repository.

\bibitem[Hurd(1989)]{hurd1989}
Hurd, M.~D. (1989).
Mortality risk and bequests.
\emph{Econometrica} 57(4), 779--813.

\bibitem[Kopczuk and Lupton(2007)]{kopczuk2007}
Kopczuk, W. and J.~P. Lupton (2007).
To leave or not to leave: The distribution of bequest motives.
\emph{Review of Economic Studies} 74(1), 207--235.

\bibitem[Kv\ae rner(2023)]{kvaerner2023}
Kv\ae rner, J. (2023).
The intergenerational transmission of wealth: Bequest motives and
child-directed transfers.
\emph{Review of Financial Studies} (Norwegian administrative-data study of
parental cancer diagnoses).

\bibitem[Lockwood(2018)]{lockwood2018}
Lockwood, L.~M. (2018).
Incidental bequests and the choice to self-insure late-life risks.
\emph{American Economic Review} 108(9), 2513--2550.

\bibitem[McGarry(1999)]{mcgarry1999}
McGarry, K. (1999).
Inter vivos transfers and intended bequests.
\emph{Journal of Public Economics} 73(3), 321--351.

\bibitem[Nakajima and Telyukova(2017)]{nakajima2017}
Nakajima, M. and I.~A. Telyukova (2017).
Reverse mortgage loans: A quantitative analysis.
\emph{Journal of Finance} 72(2), 911--950.

\bibitem[Scholz and Seshadri(working paper)]{scholzseshadri}
Scholz, J.~K. and A.~Seshadri.
Children and household wealth. Working paper.

\bibitem[Wilhelm(1996)]{wilhelm1996}
Wilhelm, M.~O. (1996).
Bequest behavior and the effect of heirs' earnings: Testing the altruistic
model of bequests.
\emph{American Economic Review} 86(4), 874--892.

\end{thebibliography}
\endgroup

\end{document}
